\section{Materials and Methods}\label{sec:methods}

This section describes, in turn, the dataset, the engagement label whose construction the audit turns on, the preprocessing and synchronisation that produce the paired epochs, the system under audit, the reference encoders and tuned baselines against which it is placed on one footing, the eight steps of the audit, and the evaluation protocol and statistics that every result shares.

\subsection{Dataset and Participants}\label{sec:data}

The NeuMa dataset \cite{ref35} contains synchronously acquired EEG and eye-tracking recordings from 42 participants (45 recruited, three excluded by the dataset authors for artefactual contamination) who browsed six digital supermarket brochure pages of 24 products each and selected, by clicking, the products they intended to buy; the sex and age distributions are those of the dataset publication. The release identifies the participants as S01 to S44 with S04 and S11 unused, and all 42 are analysed here; the raw and preprocessed releases are cited as datasets \cite{ref59,ref60}. Each brochure page is a fixed grid of product placements, so a region-of-interest (ROI) layout is defined over every stimulus; Fig.~\ref{fig1} shows one synchronised epoch of the participant used later as the HIGH-engagement exemplar, with the gaze scan-path over the page, a subset of the EEG channels and the concurrent gaze and pupil traces. The dataset characteristics used here are listed in Table~\ref{tab1}.

\begin{figure*}[!htbp]
\centering
\includegraphics[width=\textwidth]{fig1_real_overlay.pdf}
\caption{A synchronised multimodal epoch (participant S24): the brochure page with the gaze scan-path overlaid (left), a subset of the retained scalp-EEG channels (top right) and the concurrent gaze-position and pupil-diameter traces (bottom right), all within the same 5\,s epoch.}
\label{fig1}
\end{figure*}

EEG was recorded with a Wearable Sensing DSI-24 headset at 300~Hz. Of its 24 physical channels, 19 are scalp electrodes of the international 10--20 system (Fp1, Fp2, F3, F4, C3, C4, P3, P4, O1, O2, F7, F8, T3, T4, T5, T6, Fz, Cz, Pz); the remaining channels are not part of this scalp montage (three auxiliary and electrooculography channels, the A2 mastoid recorded for optional re-referencing, and a trigger channel). As in every prior analysis of this dataset \cite{ref18,ref36}, the 19 scalp electrodes are the nodes of every graph and the other channels are discarded, since the trigger channel in particular would otherwise inject stimulus-timing information into the connectivity. This selection is enforced in the dataset loader and every result reported here comes from runs under it; a run with all 24 recorded channels as nodes gives the same result within 0.01 in balanced accuracy and 0.02 in ROC-AUC, so the exclusion is a matter of principle rather than of accuracy. Eye tracking was recorded at 120~Hz as gaze position and pupil diameter.

\begin{table*}[!htbp]
\caption{NeuMa dataset characteristics as used in this study.}\label{tab1}
\small
\begin{tabular*}{\linewidth}{@{\extracolsep{\fill}}lp{8.5cm}@{}}
\toprule
Characteristic & Value \\
\midrule
Participants & 42 (release identifiers S01--S44, S04 and S11 unused; 37 evaluable under the pooled engagement index; 42 under the subject-median, positive-control and product-selection labels) \\
EEG channels & 19 (10--20 scalp electrodes) \\
EEG sampling rate & 300~Hz \\
Eye-tracking features & 3 (gaze $x$, gaze $y$, pupil diameter) \\
Eye-tracking sampling rate & 120~Hz \\
Epoch length & 5.0~s, non-overlapping \\
Engagement epochs & 385, one per brochure-page view (first 5~s after the page onset), 6--20 per participant (193 HIGH, 192 LOW) \\
Positive-control epochs & resting state versus browsing, 42 participants \\
Product-selection epochs & 2{,}338 (418 selected), one per participant--product pair \\
\botrule
\end{tabular*}
\end{table*}

\subsection{The Engagement Label and Its Construction}\label{sec:label}

Every question of this paper depends on what the engagement label is, so its provenance and construction are stated before the models. The NeuMa release annotates product selection and questionnaire responses; it contains no engagement annotation, and we are aware of no prior engagement annotation for these recordings. The index used here is therefore our own construction: a fixed rule over EEG band-power and gaze statistics, implemented in the released pipeline (\texttt{engagement\_phase3d.py}; Supplementary Table~S5) and computed on the fixed 5-s epochs of Section~\ref{sec:preproc}. Every statement about ``engagement'' in this paper refers to this rule-defined index, and the predictable properties of the index that the audit reports are properties of a constructed target, not findings about an established measure of engagement. For each epoch, five EEG terms are computed from the mean of the P3, C3, F3 and Fz channels: Welch band power of that mean signal in theta (4--8~Hz), alpha (8--13~Hz) and beta (13--30~Hz), the theta/beta ratio, and an asymmetry term equal to the difference in mean amplitude between P3 and Fz. The pipeline names these channels frontal; two of the four are parietal and central electrodes, and the asymmetry term, a time-domain mean difference of band-passed, average-referenced signals, spans $-1.05$ to $1.60$ in raw units; after min--max scaling it carries 1.5\% of the variance of the composite score, and removing it would change 18 of the 385 labels. Five gaze terms are computed from the epoch's own gaze sequence: mean absolute horizontal and vertical gaze position, horizontal gaze dispersion, the number of consecutive samples whose horizontal displacement is below 0.05 of the screen width, and the Shannon entropy of a 20-bin horizontal gaze histogram. Pupil diameter does not enter the label. The ten terms are min--max normalised over the pooled epochs of all participants and combined with fixed weights (theta $+1.5$, beta $+1.2$, alpha $-1.5$, theta/beta ratio $+1.0$, asymmetry $+0.5$, the four positive gaze terms $+1.0$ each, gaze entropy $-1.0$), and the composite is binarised at the pooled median, giving 385 epochs with 193 HIGH and 192 LOW.

Three properties of this construction shape the audit, and two of them must be kept apart. The first is label--input coupling: the label is a fixed rule over statistics that the models also receive as input, so a decoder can score well by recovering the rule from either modality. This limits the construct validity of the index as a measure of engagement, and it is why the linear recoverability of the label from its own EEG terms and from its own gaze terms is measured directly (Section~\ref{sec:audit}, step~1) and every learned score is read against the probe on the terms the model sees. The second is an evaluation dependence, separate from the coupling: because the scaler and the median are fitted on the pooled epochs, each held-out participant contributes to the construction of its own labels. Refitting both on the training participants of each fold changes 9 of the 385 labels (2.3\%), and the audit repeats the principal runs under that fold-wise label (step~4) so that this dependence is measured rather than assumed away. Third, five gaze terms and five EEG terms enter with comparable weights, but nothing in the weights says which modality carries the variance of the composite; that is an empirical question, answered in Section~\ref{res:rq1}.

Under the pooled median, five participants (S16, S31, S33, S41 and S44) carry only one class and are excluded from the test folds while remaining in every training pool, leaving 37 evaluable participants and 347 held-out epochs. As a coverage check, each participant is additionally split at their own rank-based median, which gives both classes to all 42 participants and captures within-participant relative engagement rather than an absolute level; this subject-median regime is reported among the robustness checks (Section~\ref{res:robust}). The two other labels used by the audit, the resting-state-versus-browsing contrast and the product-selection label, are defined with the audit steps that use them (Section~\ref{sec:audit}, steps~7 and~8).

\subsection{Preprocessing and Synchronisation}\label{sec:preproc}

EEG was re-referenced to the common average, band-pass filtered to 1--45~Hz and notch filtered at 50~Hz; ocular and muscular artefacts were attenuated by independent component analysis with 15 components. For one participant (S01) four non-cortical channels flagged as noisy were zero-filled; no scalp electrode was removed for any participant. Signals were kept at 300~Hz. Per-channel $z$-score normalisation was applied within each participant from that participant's own statistics, without reference to the labels; because this uses the statistics of the whole held-out recording, the audited decoder was also retrained with per-epoch instance normalisation, which uses no participant-level statistic and reaches the same accuracy (Section~\ref{res:robust}). Raw gaze coordinates and pupil diameters were cleaned by discarding samples lost to blinks and tracking drop-outs, with short gaps linearly interpolated; guided by a feature-ablation analysis, the un-normalised three-channel sequence (gaze $x$, gaze $y$, pupil) was retained as the eye-tracking input, since higher-dimensional configurations did not improve subject-independent performance.

The two streams were aligned on shared stimulus-onset markers and kept at their native rates. Engagement epochs are cut at the brochure-page-onset markers: each page view yields one epoch spanning the first 5.0~s after the marker, carrying paired EEG ($19\times1{,}500$) and eye-tracking ($3\times600$) tensors, and a participant contributes as many epochs as page views, revisits included. The 42 participants viewed the six pages 385 times in total (91, 60, 58, 62, 57 and 57 views of pages 1 to 6), between 6 and 20 views per participant and about nine on average, so the engagement analyses use about 46~s of recording per participant; the remainder of each page view enters only the product-selection analysis of audit step~8, which uses the full views. Fixed windows were preferred to the gaze-defined, variable-duration segments of prior NeuMa work \cite{ref18} for two reasons: they yield equal-shape tensors for batched graph construction and transformer encoding, and, because the label is derived in part from gaze statistics, a gaze-defined duration would tie epoch length to the label, whereas a fixed window removes that dependence. Two derived representations are used by several models. The per-window band-power node features are the relative power of each electrode in five canonical bands (each band's share of the window's total power, layer-normalised per window) over the ten 0.5-s windows of an epoch; they are the node features of the decoder's graphs and the input of the linear encoder used in the positive control. The ROI saliency vector $r\in\mathbb{R}^{10}$ is the normalised count of gaze samples falling in each cell of a fixed $5\times2$ grid over the page, a gaze-occupancy histogram rather than a fixation-based dwell time; the 24 product boxes drawn in Fig.~\ref{fig1} are the brochure's layout and not a model input.

\subsection{The System Under Audit}\label{sec:system}

INS-HDGS-CMT (Interpretable Neuro-Symbolic Hybrid Dynamic Graph Spiking Cross-Modal Transformer) is an end-to-end multimodal architecture of six components (Fig.~\ref{fig2}). It is described as the object of study: each component is introduced with the purpose it was built for and the audit step that tests that purpose, and Table~\ref{tab_pipeline} records the outcome beside each stage so that the Methods do not describe the architecture as more than the Results support. The full specification, with every equation, is given in Supplementary Note~S1; the equations retained here are those the Results refer to. Notation is summarised in Supplementary Table~S4 and the hyperparameters, fixed before the reported runs and never searched per fold, in Table~\ref{tab2}.

\begin{figure*}[!htbp]
\centering
\includegraphics[width=\textwidth]{fig2_Architecture.pdf}
\caption{Overview of the system under audit. Signal preprocessing, synchronisation and epoch construction (A) yield paired EEG and eye-tracking tensors that feed the EEG graph--spiking pathway (B) and the eye-tracking attention pathway (C); these are fused by the cross-modal transformer (D), and the neuro-symbolic decision layer (E) yields the engagement prediction (F), one per 5-s epoch.}
\label{fig2}
\end{figure*}

\begin{table*}[!htbp]
\caption{Processing pipeline of the system under audit: operation, design purpose and the audit outcome for each stage.}\label{tab_pipeline}
\footnotesize
\begin{tabular}{@{}p{0.3cm}p{2.9cm}p{5.2cm}p{5.2cm}@{}}
\toprule
\# & Stage & Operation & Designed purpose and audit outcome \\
\midrule
1 & Dynamic EEG graph construction & Per-window Pearson correlation between electrode signals $\rightarrow$ sequence of graphs $\{G_t\}$, Eqs.~\eqref{eq2pearson}--\eqref{eq3}. & Time-varying connectivity rather than a static matrix. Ablating the pathway costs 0.09 in balanced accuracy; density-matched static or random topologies produce no detectable reduction, so the measured topology was not necessary for that accuracy (Section~\ref{res:rq3}). \\
2 & Eye-tracking encoding & Gaze-$x$, gaze-$y$ and pupil streams $\rightarrow$ attention tokens; the ROI saliency vector gates the EEG embedding and modulates the graph adjacency. & Overt visual attention. Accounts for the accuracy on this label: removing the gaze inputs lowers ROC-AUC by 0.29 (Section~\ref{res:rq1}). \\
3 & Spiking front-end (LIF) & LIF layers over the ten per-window feature vectors $\rightarrow$ event-driven embedding merged with the graph embedding, Eq.~\eqref{eq:lif}. & Event-driven representation. Accuracy unchanged by its removal; 0.55 ROC-AUC on its own on the positive control; slower than a dense twin on the hardware used (Section~\ref{res:rq3}). \\
4 & Hierarchical graph attention & Attention over electrodes and edges across the graph sequence. & Per-window electrode and edge weighting with inspectable coefficients. The coefficients are descriptive; the topology nulls show the decision does not depend on them (Section~\ref{res:rq3}). \\
5 & Temporal transformer & Self-attention over the sub-window sequence. & Temporal evolution of connectivity across the epoch. \\
6 & Cross-modal fusion & Self-attention over EEG, graph and eye-tracking tokens, then EEG-to-context cross-attention and gated residual fusion, Eq.~\eqref{eq:fusion}. & Lets the EEG representation query gaze context. No detectable difference from bidirectional cross-attention or from a gaze-only encoder after correction (Sections~\ref{res:rq1} and~\ref{res:rq3}). \\
7 & Neuro-symbolic decision layer & Soft match to eight rule premises; rules vote HIGH/LOW; gated aggregation with a bypass classifier, Eqs.~\eqref{eq17}--\eqref{eq18}. & Inspectable rule-activation traces. In the trained model the decision is bypass-dominated; attribution summaries are reported from the variant trained with the gate closed, at a cost of 0.04 in ROC-AUC (Sections~\ref{res:rq3} and~\ref{res:rules}). \\
8 & Engagement prediction & Logits $\rightarrow P(\mathrm{HIGH})$; threshold transferred from a validation participant, never from the test participant. & Subject-independent output (Section~\ref{sec:eval}). \\
\botrule
\end{tabular}
\end{table*}

\paragraph{Dynamic functional graphs.} Each 5-s epoch is divided into ten non-overlapping 0.5-s windows (150 samples). In window $t$ the EEG is represented as a graph $G_t=(V,E_t,A_t)$ whose node set $V$ is the 19 electrodes and whose edges reconfigure over time. Edge weights are the Pearson correlation between the windowed time-domain signals of each electrode pair,
\begin{equation}
\rho_{ij}(t) = \frac{\operatorname{Cov}\!\left(x_i(t), x_j(t)\right)}{\sigma_{x_i(t)}\,\sigma_{x_j(t)}}, \label{eq2pearson}
\end{equation}
computed independently per window; the phase-locking value is an alternative measure that was evaluated at reduced scale (Section~\ref{res:robust}; Supplementary Note~S1). The weighted matrix $C_t(i,j)=\rho_{ij}(t)$ is thresholded to a binary adjacency,
\begin{equation}
A_t(i,j) = \begin{cases} 1, & |C_t(i,j)| \geq \tau \text{ or } i=j, \\ 0, & \text{otherwise}, \end{cases} \label{eq3}
\end{equation}
with $\tau=0.30$ fixed before the reported runs for every participant and fold; the graph attention consumes $A_t$ as a neighbourhood mask, so the correlation magnitude decides whether an edge exists but not the weight passed along it, which the attention learns. Over the 3{,}850 windows of the 385 epochs this threshold retains 57\% of the possible edges on average, and its structural and downstream consequences are examined in Section~\ref{res:robust}. The node features are the per-window band-power vectors of Section~\ref{sec:preproc}. Fig.~\ref{fig3} illustrates the construction on one real epoch. Treating the topology itself as time-varying, rather than building one graph per trial as related graph-based EEG models do \cite{ref19,ref20,ref21,ref22}, was the design premise of the pathway; whether the measured topology is what the pathway uses is the question that the density-matched nulls of audit step~5 answer.

\begin{figure*}[!htbp]
\centering
\includegraphics[width=\textwidth]{fig_graph.pdf}
\caption{Dynamic functional-graph construction from one real epoch (Pearson connectivity, Eq.~\eqref{eq2pearson}). (A) Ten 0.5~s windows of a 5~s epoch; the shaded window is shown in B--D. (B) Two $z$-scored electrode signals, whose correlation gives one edge weight. (C) The connectivity matrix $|C_t|$ of the window. (D) The graph $G_t$ after thresholding at $\tau=0.30$; node size scales with degree.}
\label{fig3}
\end{figure*}

\paragraph{Spiking front-end.} In parallel with the graph attention, a two-layer leaky integrate-and-fire (LIF) stack \cite{ref38} of 128 units receives a learned band-weighted sum of the per-window node features, one 19-electrode vector per window, and is run for ten time steps. The discrete-time dynamics are
\begin{equation}
U_t = \lambda V_{t-1} + I_t, \qquad S_t = \mathbf{1}\!\left[U_t \geq V_{\mathrm{th}}\right], \qquad V_t = U_t\left(1 - S_t\right), \label{eq:lif}
\end{equation}
with pre-reset potential $U_t$, membrane leak $\lambda\in(0,1)$, a spike when $U_t$ crosses the threshold, a hard reset to zero after a spike, and a surrogate gradient through the spike function \cite{ref39}; the spike trains are layer-normalised between layers and the resulting embedding is concatenated with the graph embedding to form the EEG representation. The front-end was included for an event-driven representation with an inductive bias toward salient transients \cite{ref40,ref41,ref42}. Its effect within the decoder is measured by ablation, its cost on conventional hardware by direct timing, and its capacity on its own by a separate configuration applied to the raw 1{,}500-sample epoch pooled to 300 steps, which is used only in the positive control (audit steps~5 and~7).

\paragraph{Graph attention and temporal transformer.} The graphs are processed by graph attention \cite{ref43} in two layers (four heads, then one), which weight each node's neighbours by learned coefficients computed from the projected node features and normalised over the neighbourhood defined by $A_t$; the coefficients double as a descriptive importance signal. The sequence of graph-level embeddings is then passed through a transformer encoder \cite{ref44} with four heads and a 128-dimensional embedding, so that every pair of sub-windows can interact directly. Both operations are specified in full in Supplementary Note~S1. Graph attention was chosen over spectral graph convolution because it accommodates the per-window adjacency changes of Eq.~\eqref{eq3} without retraining filters; the internal graph-attention baseline that applies the same mechanism to a single static graph per epoch is among the baselines of Section~\ref{sec:references}.

\paragraph{Eye-tracking pathway.} The eye-tracking encoder consumes the raw three-channel sequence (600 samples) and produces attention tokens; it is not itself gated by the ROI vector. The ROI saliency vector $r$ enters the EEG side: it gates the merged EEG embedding through a learned sigmoid gate, $\tilde h = h \odot \sigma(W_r r + b_r)$, and modulates the graph adjacency before attention, which is the second route by which gaze information enters the neural pathway (Fig.~\ref{fig2}); the eye-tracking encoder additionally predicts a soft ROI distribution that is used only by the inactive ROI-guidance term. Fixation duration and dwell time, the behavioural indices used in the descriptive and case-study analyses, are defined in Supplementary Note~S1; they are not model inputs. The ROI pathway is active only in the full multimodal configuration and is disabled entirely in the gaze-free variant of audit step~2.

\paragraph{Cross-modal fusion.} The EEG, graph and eye-tracking embeddings are treated as a three-token sequence and processed by multi-head self-attention, through which the modalities exchange information mutually, and then by a directed cross-attention in which the EEG representation queries the graph and eye-tracking context,
\begin{align}
Z_{\mathrm{EEG}\to m} &= \mathrm{Attention}\!\left(Q_{\mathrm{EEG}},K_{m},V_{m}\right), \quad m\in\{\mathrm{graph},\mathrm{ET}\}, \nonumber\\
Z_{\mathrm{fusion}} &= \left[Z_{\mathrm{EEG}} \,;\, Z_{\mathrm{EEG}\to\mathrm{graph}} \,;\, Z_{\mathrm{EEG}\to\mathrm{ET}}\right], \label{eq:fusion}
\end{align}
followed by a learnable gated residual that passes modality information selectively (Fig.~\ref{fig4}). Anchoring the fused embedding on the EEG stream is what lets the gaze-free variant be obtained by removing the eye-tracking input alone. The fusion was designed to capture the conditional dependence between gaze and neural activity that concatenation cannot; whether it sets the level of accuracy is tested against simpler fusion designs and against a gaze-only encoder (audit step~3).

\begin{figure*}[!htbp]
\centering
\includegraphics[width=\textwidth]{fig4_cross-modal.pdf}
\caption{Cross-modal fusion and neuro-symbolic decision layer of the system under audit: self-attention over the EEG, graph and eye-tracking tokens, directed cross-attention with the EEG representation as query, gated residual fusion, and the eight-rule module; the learned bypass blend of Eq.~(6) acts at the output and is not drawn. In the trained model the decision is bypass-dominated (Section~\ref{res:rq3}).}
\label{fig4}
\end{figure*}

\paragraph{Neuro-symbolic decision layer.} The fused representation $z\in\mathbb{R}^D$ is projected to a rule key $k=\mathrm{Dropout}(\mathrm{GELU}(W_k z + b_k))$ and matched to $n_r=8$ learnable rule premises $q_r$ by a softmax over the scaled similarities $s_r=q_r^{\top}k/\sqrt{D}$, giving activations $a_r\in[0,1]$ that sum to one. Each rule emits a class-level conclusion $\ell_r=W_r k+b_r$ through its own linear head, and the rule evidence is their activation-weighted sum,
\begin{equation}
R = \sum_{r=1}^{n_r} a_r \ell_r . \label{eq17}
\end{equation}
A bypass classifier on $z$ is blended with the rule evidence through a learnable gate,
\begin{equation}
\hat y = \mathrm{softmax}\!\left(\alpha(W_b z + b_b) + (1-\alpha) R\right), \qquad \alpha=\sigma(\alpha_0), \label{eq18}
\end{equation}
initialised at $\alpha=0.57$; as $\alpha\to0$ the decision becomes rule-driven and as $\alpha\to1$ it reverts to the unconstrained network. Diversity and sparsity regularisers are intended to keep the premises distinguishable and the activations sparse; in the gated model the activations remained uniform (Section~\ref{res:rq3}; Supplementary Note~S1). At inference each rule's activation is attributed to the physical inputs by integrated gradients \cite{ref54} (zero baseline, 16 steps): the largest-magnitude entries of the resulting electrode-by-band map, with their sign, are reported as an attribution summary of rule $r$, and, as its conclusion, the class on which its mean activation is higher over the held-out epochs; the summary describes the inputs to which the activation is most sensitive and is not a validated logical premise. The module was designed to replace an opaque classifier with inspectable rules, in the spirit of interpretable-by-design and neuro-symbolic models \cite{ref45,ref46,ref47}; because the gate could in principle remove the rules from the decision, its learned value, the agreement between $\arg\max R$ and the gated decision, the accuracy with $\alpha$ forced to zero, and a variant trained throughout with $\alpha\equiv0$ are all measured (audit step~6).

\paragraph{Objectives and training.} The classification loss is a focal, class-balanced cross-entropy \cite{ref48,ref49} with label smoothing 0.05, chosen because the per-fold training pool and the small held-out set are often skewed even though the pooled data are balanced. Three auxiliary terms shape the representation: a maximum mean discrepancy (MMD) penalty between the EEG embeddings of different participants within a mini-batch and a domain-adversarial participant classifier with gradient reversal, both aimed at the across-participant shift that subject-independent decoding must survive \cite{ref31,ref32}, and a cross-modal contrastive term (InfoNCE, temperature 0.07, hard-negative weighting) that makes the EEG embedding of an epoch predictive of its own eye-tracking embedding, which is the training-time route by which gaze shapes the neural pathway. The total objective is
\begin{multline}
L_{\mathrm{total}} = L_{\mathrm{cls}} + \lambda_{\mathrm{rules}} L_{\mathrm{rules}} + \lambda_{\mathrm{con}} L_{\mathrm{con}} \\
{} + \lambda_{\mathrm{mmd}} L_{\mathrm{mmd}} + \lambda_{\mathrm{adv}} L_{\mathrm{adv}}, \label{eq26}
\end{multline}
with weights fixed before the reported runs ($\lambda_{\mathrm{rules}}=0.02$, $\lambda_{\mathrm{con}}=0.05$, $\lambda_{\mathrm{mmd}}=\lambda_{\mathrm{adv}}=0.10$) and the contrastive and rule terms ramped in over the first 40\% of training; an ROI-guidance term was implemented but is evaluated only at inference and contributes no gradient. The MMD term is evaluated only when both sides of a comparison contain at least eight epochs, which at batch size 32 is met in few mini-batches, so it acts sparsely. Training uses AdamW with cosine annealing, gradient clipping, early stopping on a validation participant's balanced accuracy, deterministic half-epoch circular-shift augmentation with stochastic amplitude scaling, noise, channel drop-out and temporal masking, and a five-member seed-averaged ensemble whose probabilities are averaged (Table~\ref{tab2}). The model has 1.72 million parameters; single-epoch inference on a commodity CPU (AMD Ryzen~7, no GPU) takes $14.9\pm1.3$~ms, and batching reaches 156 epochs per second, so the decoder runs in real time without GPU acceleration. Experiments were implemented in Python 3.12 with PyTorch 2.7.1 and CUDA 12.6.

\begin{table*}[!htbp]
\caption{Implementation details of the system under audit: architecture, optimiser and training hyperparameters held fixed across all folds.}\label{tab2}
\small
\begin{tabular*}{\linewidth}{@{\extracolsep{\fill}}lp{0.68\linewidth}@{}}
\toprule
Component & Configuration \\
\midrule
Epoch length & 5.0~s \\
Batch size & 32 \\
Learning rate & $5\times10^{-4}$ \\
Optimiser & AdamW \\
Training epochs & at most 250 \\
Dropout & 0.25 (0.17 in the eye-tracking encoder, 0.08 in the ROI, fusion and rule blocks) \\
Weight decay & 0 \\
Schedule & Cosine annealing with warm restarts ($T_0=50$ epochs, $\eta_{\min}=10^{-6}$); gradient-norm clipping at 1.0 \\
Early stopping & Validation balanced accuracy, patience 50, minimum 20 epochs \\
Transformer heads & 4 \\
Hidden / embedding dim & 128 \\
GAT heads (layer 1 / 2) & 4 / 1 \\
Connectivity threshold $\tau$ & 0.30 (sensitivity: Section~\ref{res:robust}) \\
ROI saliency grid & $5\times2$ (10 cells; sensitivity: Section~\ref{res:robust}) \\
Rule bypass gate $\alpha$ & learned (sigmoid, initialised at 0.57); fidelity: Section~\ref{res:rq3}, Supplementary Table~S13 \\
Loss function & Focal cross-entropy, $\gamma=3$, class-balance factor 1, label smoothing 0.05, effective-number class weights \cite{ref49} \\
Random seed (base) & 42 \\
Ensemble members & 5 (seed-averaged) \\
\botrule
\end{tabular*}
\end{table*}

\subsection{Reference Encoders and Tuned Baselines}\label{sec:references}

Three standard encoders anchor the findings. EEGNet \cite{ref50} and ShallowConvNet \cite{ref51} are the compact convolutional EEG encoders most widely used as references in EEG decoding, and they bracket the EEG family here: by average rank on balanced accuracy, EEGNet places last and ShallowConvNet first among the EEG encoders on the engagement index (Section~\ref{res:rq2}). The eye-tracking LSTM (ET-LSTM) is a single-layer long short-term memory network over the raw three-channel gaze sequence with a linear head, the eye-tracking encoder with the highest ROC-AUC on this index. To show what a multimodal decoder reaches with off-the-shelf parts, each convolutional encoder was combined with the ET-LSTM in a decision-level late fusion that averages the two encoders' logits (EEGNet+ET-LSTM and ShallowConvNet+ET-LSTM), with both encoders trained jointly under the same tuning rule as every other baseline. These fusions have no cross-modal attention, no graph, no spiking front-end and no rule module, so the difference between them and the decoder is the whole of its architecture.

The remaining baselines make the comparison a field rather than a pair. Eight EEG encoders receive no gaze input: EEGNet, ShallowConvNet, DeepConvNet \cite{ref51}, CNN-LSTM, CNN-BiLSTM, an EEG Transformer, TSception \cite{ref52} and a graph attention network on a single static graph per epoch (GAT); BrainGCN, at chance on this label (balanced accuracy 0.47, ROC-AUC 0.47), enters the Holm-corrected comparisons but not the nine-encoder Nemenyi family. Three eye-tracking encoders (ET-LSTM, ET-GRU, ET-Transformer) receive gaze alone. Six fusion architectures receive both modalities: cross-attention fusion, a multimodal transformer, a dynamic-GAT with an eye-tracking transformer, a dual transformer, an early-fusion multilayer perceptron and a late fusion of CNN-LSTM with ET-LSTM. With the two new fusions the baseline set has twenty members.

Every baseline was tuned per fold with matched effort: twelve hyperparameter configurations (learning rate, weight decay, dropout and batch size, from an architecture-specific space that always includes the common budget of Adam, $10^{-3}$, $10^{-4}$, 32) were each trained as a single model with class-weighted cross-entropy and early stopping on the validation participant's balanced accuracy (patience 30, at most 250 epochs), and the configuration with the highest validation balanced accuracy was applied to the test participant. The decoder's hyperparameters were fixed before the reported runs and not searched (Table~\ref{tab2}); it additionally uses the five-member ensemble, the focal loss and the augmentation of Section~\ref{sec:system}, which the baselines do not, so the protocols are matched in effort rather than identical, and this asymmetry is stated where the comparisons are read. Selection never touched the test participant, and the selected configurations are reported per fold in the repository. Results under a common fixed budget without per-fold tuning are given in Supplementary Table~S8.

\subsection{The Audit Protocol}\label{sec:audit}

The audit consists of eight controls. Each is stated with the question it answers and the reading its outcomes admit, so that the Results can report outcomes without re-deriving their logic. None depends on a property of the audited architecture; every step applies to any decoder that takes the same inputs, and all are implemented in the released code.

\begin{enumerate}
\item \textbf{Label recoverability per modality} (RQ1). Logistic-regression probes are trained under LOSOCV on the label's own five EEG terms and, separately, on its own five gaze terms, computed within each epoch (\texttt{label\_leakage\_audit.py}). Their ROC-AUC is the reference for each modality. A learned model that does not exceed the probe on the terms it sees has not been shown to use information beyond them; one that exceeds it may use additional information or may use the same terms non-linearly. The comparison orders performance and does not identify which information a model uses.
\item \textbf{Gaze-free variant} (RQ2). The decoder is retrained with every gaze-derived input removed: no eye-tracking sequence, no ROI vector, no contrastive term (which needs the gaze embedding) and no MMD term, with the adversarial term retained. It is the only configuration of the architecture independent of the label's gaze terms, and its accuracy is the EEG-only capacity of the design. Because the MMD term is removed with the gaze inputs, the MMD ablation of step~5 bounds that confound.
\item \textbf{Reference encoders on one footing} (RQ1, RQ2). Every model of Section~\ref{sec:references}, the decoder and its gaze-free variant are evaluated on identical folds, labels and operating points, and compared by paired tests with Holm correction within each family, by a Friedman test with a Nemenyi critical-difference analysis over the EEG family, and by a per-participant scatter of the principal models.
\item \textbf{Fold-wise relabelling} (RQ2). The label's scaler and median are refitted on the training participants of each fold, and the validation participant is drawn by a rule that is identical for every model and never selects a single-class participant. The decoder, its gaze-free variant, EEGNet, ShallowConvNet and ET-LSTM are retrained under this protocol and compared with their pooled-label runs by paired differences with bootstrap intervals; the decoder's run-to-run spread (Section~\ref{sec:eval}) is reported beside its own differences as a reference.
\item \textbf{Component ablation and topology nulls} (RQ3). Each component is removed in turn and the model retrained on the same folds. For the graph pathway, two density-matched nulls change only the adjacency: a static graph that keeps the $k$ most persistent edges of the epoch, with $k$ the mean edge count per window, repeated over all windows, and, per window, a uniform random graph with exactly the same number of edges (Erd\H{o}s--R\'enyi $G(n,m)$) onto which the epoch's edge weights are re-dealt. Node features, ROI vector and gaze are untouched, so the nulls test whether the measured time-varying topology is what the pathway uses. For the spiking front-end, wall-clock latency is measured against a dense twin of identical layer widths on the CPU and on an A100 GPU (\texttt{snn\_energy\_measured.py}).
\item \textbf{Rule fidelity} (RQ3). Four measurements are taken on the rule module: the learned gate value $\alpha$ and the distribution of activations $a_r$; the agreement between the rule-only decision $\arg\max R$ and the gated decision; the accuracy when $\alpha$ is forced to zero at inference; and the accuracy and rule traces of a variant trained throughout with $\alpha\equiv0$ (\texttt{rule\_fidelity.py}). Attribution summaries are reported only from a variant whose decision passes through the rule module.
\item \textbf{Positive control} (RQ2). On the same recordings, a contrast the EEG is expected to carry is decoded: the two-minute pre-brochure resting-state fixation versus brochure browsing, in 5-s epochs, 42 participants, LOSOCV. A ladder of encoders is trained gaze-free: a linear encoder on the band-power node features, the graph pathway with a plain head, the spiking front-end alone on raw EEG, the spiking front-end with the graph pathway, the gaze-free variant with its MMD and adversarial regularisers retained, EEGNet and ShallowConvNet, and the full decoder; ensemble sizes differ by rung and are given with the results (Supplementary Table~S16). The control tests whether the recordings and the pipelines carry condition-discriminating information at all. It does not isolate neural engagement: resting fixation and browsing differ in visual input, eye movements and muscle activity as well as in cortical state, so a positive result shows that usable information is present, not what it is.
\item \textbf{Product-selection label} (RQ4). Whether the viewed product was selected for purchase in the simulated shopping task (NeuMa questionnaire item Q77, which agrees with the recorded selection clicks for 99\% of pairs; the task records intended purchase, not a verified transaction) defines a label with one 5-s epoch per participant--product pair (\texttt{product\_epoching.py}; 2{,}338 epochs, 418 selected, 42 participants). Gaze samples are assigned to the product boxes of each page, contiguous on-product runs are formed (gaps up to 100~ms bridged, runs shorter than 0.3~s discarded; products with less than 1~s of total dwell, the summed run duration over the session, are skipped), and the epoch is anchored 1~s before the onset of the product's longest eligible run on the page, with the same anchor formula for both classes; for selected products only runs whose window ends at least 250~ms before the first click on that product are eligible, so no click on that product falls inside the epoch (preparatory or other motor activity is not excluded), and for unselected products the longest run is used, so eligibility is class-dependent. The anchor is chosen over the whole page view and is therefore retrospective, while the epoch content precedes the selection click. The decoder, EEGNet, ShallowConvNet, ET-LSTM, EEGNet+ET-LSTM and ShallowConvNet+ET-LSTM are trained on it under LOSOCV, and compared with two families of logistic probes: within-epoch probes on the five EEG and five gaze statistics of the engagement rule, which are not terms of this label and serve only as within-epoch reference features, and session-level probes on how long the product was looked at over the whole session, alone and with the number of gaze runs, views and the anchor run. The session-level dwell accumulates over the session, including viewing after the selection, so the session-level comparison is an association between looking time and selection, not a prospective prediction from a common cutoff. An incremental-validity probe adds each model's held-out epoch probability to the dwell probe in a second LOSOCV logistic regression; because those probabilities come from models whose training folds included the other participants of each second-level training set, the stacking is not nested in the outer participant split; the dependence is expected, if anything, to favour the model, but the increment is not a leakage-free estimate.
\end{enumerate}

Steps~1 to~4 attribute the accuracy to a modality, steps~5 and~6 to a component, and steps~7 and~8 place the index against a contrast the recordings do carry and against a behavioural criterion. Section~\ref{sec:intro} lists the question that each step answers, and Section~\ref{sec:results} reports the steps in that order.

\subsection{Evaluation and Statistics}\label{sec:eval}

All models are evaluated under LOSOCV: every epoch of one held-out participant forms the test set and the epochs of all other participants the training pool, so no held-out epoch is used for parameter fitting or hyperparameter selection. This is the estimate of generalisation to an unseen individual that prospective deployment requires, and it removes the inflation that subject-mixed splits produce \cite{ref13,ref28,ref29,ref30,ref62}; the residual dependence introduced by the pooled label is a separate question that audit step~4 measures. Performance is summarised as the mean $\pm$ SD over folds of accuracy, balanced accuracy, macro-averaged F1, ROC-AUC, the area under the precision--recall curve and the Matthews correlation coefficient (MCC); ROC-AUC, which is threshold-invariant, is the reporting metric for ranking quality, and balanced accuracy and MCC are reported at two operating points, neither of which uses a test label. At the first, tabulated throughout, the probability obtained from the ensemble-averaged logits is thresholded at the Youden-$J$ value found on one validation participant drawn from the training pool (the transferred threshold, 0.05--0.73 across folds); the rule falls back to 0.5 when the validation participant has fewer than four epochs or one class, which occurred in 2 of the 37 folds. At the second, the calibrated operating point, the ensemble-averaged logits are divided by a temperature fitted by maximum likelihood on the same participant's averaged validation logits, and the decision threshold is re-tuned by Youden's $J$ on the calibrated validation probabilities only when the validation participant has at least twenty epochs, otherwise it is fixed at 0.5; a positive temperature leaves the 0.5 decision unchanged, so the calibrated point coincides with a fixed 0.5 cut-off in every fold except the one whose validation participant had twenty epochs. Under the pooled-label protocol the validation participant is drawn independently for the decoder and for each baseline, so those comparisons share folds and test participants but not validation participants; under the fold-wise protocol of audit step~4 the draw is deterministic from the seed and the test participant and identical for every model, so the decoder, its gaze-free variant, EEGNet, ShallowConvNet and ET-LSTM are matched in validation participants there. Label-free per-participant thresholds, transductive and causal, are reported as a threshold-adaptation analysis (Section~\ref{res:robust}). Expected calibration error is computed over ten equal-width bins on the 347 held-out epochs pooled, because a participant contributes about nine epochs, and the Brier score beside it; both come from the saved held-out probabilities of every model through one script (\texttt{pooled\_calibration.py}).

Statistical conclusions rest on paired non-parametric tests over folds. Paired comparisons use the Wilcoxon signed-rank test with zero differences discarded, Holm--Bonferroni correction within each family, the matched-pairs rank-biserial correlation as effect size, and 95\% bootstrap confidence intervals of the mean paired difference (10{,}000 resamples); family-wise comparisons use a Friedman test with a Nemenyi critical-difference analysis (critical difference 1.97 ranks for nine encoders over 37 folds); the association between held-out predictions and labels is tested by a within-fold label permutation (5{,}000 shuffles). Two references calibrate the size of every effect. The run-to-run spread of the decoder, measured by three trainings of the same configuration with the same seed, is 0.03 in balanced accuracy and 0.02 in ROC-AUC; it is reported beside the decoder's own variants as a reference and is not used as a significance threshold. Cross-validation on a cohort of this size carries wide error bars \cite{ref63}, and the per-fold test set is about nine epochs, so single-fold and hard-participant metrics are indicative rather than precise. Differences that are not significant are reported as differences that were not detected, not as equivalence, since no equivalence margin was pre-specified and the intervals are wide. Overfitting is guarded against by the augmentation, dropout, subject-alignment terms and ensemble of Section~\ref{sec:system}, by early stopping on a held-out participant, by the permutation test, and by a participant-count learning curve in which LOSOCV is repeated on random subsets of 10 to 37 participants (Section~\ref{res:robust}). The complete pipeline, hyperparameters, per-fold selections, saved held-out probabilities and analysis scripts are released with the code (Declarations).
